\section{Conclusion}
\label{conclusion}
This report introduces \textbf{MiniOneRec}, the first fully open-source generative recommendation framework, which to the best of our knowledge provides an end-to-end workflow spanning SID construction, SFT, and recommendation-oriented RL.
By systematically validating scaling laws on public benchmarks, we demonstrate that larger generative recommenders achieve lower training and evaluation losses than their smaller counterparts, confirming the parameter–efficiency advantage of the SID-based paradigm over traditional embedding-centric models.  
Building on this insight, we introduce post-training techniques:  
(i) full-process SID alignment, which embeds SID tokens into the model vocabulary and imposes auxiliary alignment tasks across both SFT and RL stages, and (ii) reinforced preference optimization, which combines constrained decoding, beam-based sampling, and hybrid reward design.  
Extensive experiments on Amazon Review show that MiniOneRec consistently surpasses strong sequential, generative, and LLM-based baselines, while maintaining a lean post-training footprint.

Looking forward, we will keep maintaining and extending the MiniOneRec codebase.  
A public roadmap will guide future developments, and we warmly welcome community contributions.  
Planned updates include new datasets, more advanced tokenisation schemes, larger backbone models, and enhanced training pipelines, ensuring that MiniOneRec remains a solid reference platform for research and practice in large-scale generative recommendation.
